\TieudeT{PHONG TỎA VẬT LÍ 12 - VẬT LÍ NHIỆT}
\subsubsection*{PHẦN I. Thí sinh trả lời câu hỏi từ câu 1 đến câu 18. Mỗi câu hỏi thí sinh chỉ chọn một phương án}
\setcounter{ex}{0}
\begin{ex}%Câu 1
	Phát biểu nào dưới đây là đúng khi nói về những đặc điểm của chất rắn?
	\choice
	{Có khối lượng, hình dạng xác định, không có thể tích xác định}
	{Có khối lượng xác định, hình dạng và thể tích không xác định}
	{\True Có khối lượng, hình dạng, thể tích xác định}
	{Có khối lượng và thể tích xác định, hình dạng không xác định}
	\loigiai{
		
	}
\end{ex}

\begin{ex}%Câu 2
	Chất lỏng không có hình dạng xác định vì các phân tử chất lỏng
	\choice
	{dao động tại các vị trí cân bằng xác định}
	{có thể chuyển động phân tán ra xa nhau}
	{\True dao động quanh các vị trí cân bằng có thể dịch chuyển được}
	{có thể chuyển động tự do}
	\loigiai{
		
	}
\end{ex}

\begin{ex}%Câu 3
	Sự nóng chảy ở chất rắn kết tinh bắt đầu xảy ra khi
	\choice
	{một số phân tử dao động mạnh hơn các phân tử xung quanh}
	{một số phân tử va chạm với các phân tử xung quanh}
	{một số phân tử dao động mạnh lên và truyền năng lượng dao động cho các phân tử khác}
	{\True một số phân tử thắng được lực liên kết với các phân tử xung quanh và thoát khỏi liên kết với chúng}
	\loigiai{
		
	}
\end{ex}

\begin{ex}%Câu 4
	Phát biểu nào dưới đây là \textbf{không} đúng khi nói về sự nóng chảy của các chất rắn?
	\choice
	{Mỗi chất rắn kết tinh nóng chảy ở một nhiệt độ xác định không đổi ứng với một áp suất bên ngoài xác định}
	{Nhiệt độ nóng chảy của chất rắn kết tinh phụ thuộc áp suất bên ngoài}
	{Chất rắn kết tinh nóng chảy và đông đặc ở cùng một nhiệt độ xác định không đổi}
	{\True Chất rắn vô định hình cũng nóng chảy ở một nhiệt độ xác định không đổi}
	\loigiai{
		
	}
\end{ex}

\begin{ex}%Câu 5
	Một số phân tử ở gần mặt thoáng chất lỏng, chuyển động hướng ra ngoài, có...(1)... đủ lớn thắng được lực liên kết giữa các phân tử thì có thể thoát ra ngoài khối chất lỏng. Như vậy, có thể nói sự bay hơi là sự hóa hơi xảy ra ở ...(2)... của khối chất lỏng. Điền vào chỗ trống các cụm từ thích hợp.
	\choice
	{\True (1) động năng; (2) mặt thoáng}
	{(1) thế năng; (2) mặt thoáng}
	{(1) động năng; (2) trong lòng}
	{(1) thế năng; (2) trong lòng}
	\loigiai{
		
	}
\end{ex}

\begin{ex}%Câu 6
	Sự sôi của chất là
	\choice
	{sự hóa hơi của các chất ở mọi nhiệt độ}
	{sự hóa hơi của chất xảy ra trên bề mặt chất lỏng ở mọi nhiệt độ}
	{\True sự hóa hơi của chất xảy ra cả trên bề mặt và trong lòng khối chất lỏng ở một nhiệt độ xác định}
	{sự hóa hơi của chất xảy ra cả trên bề mặt và trong lòng khối chất lỏng ở mọi nhiệt độ}
	\loigiai{
	}
\end{ex}

\begin{ex}%Câu 7
	Ở nhiệt độ không tuyệt đối, tất cả các chất đều có động năng chuyển động nhiệt của các phân tử
	\choice
	{khác không và thế năng của chúng là tối thiểu}
	{bằng không và thế năng của chúng là tối thiểu}
	{bằng không và thế năng của chúng cũng bằng không}
	{\True khác không và thế năng của chúng bằng không}
	\loigiai{
		
	}
\end{ex}

\begin{ex}%Câu 8
	Trong quá trình thực hiện công, ...(1)... từ dạng năng lượng khác (ví dụ cơ năng chẳng hạn) sang ...(2)... Điền vào chỗ trống cụm từ thích hợp.
	\choice
	{ (1) không có sự chuyển hóa; (2) nội năng}
	{ (1) có sự chuyển hóa; (2) hóa năng}
	{ (1) không có sự chuyển hóa; (2) hóa năng}
	{\True (1) có sự chuyển hóa; (2) nội năng}
	\loigiai{
		
	}
\end{ex}

\begin{ex}%Câu 9
	Dùng tay nén pít-tông đồng thời nung nóng khí trong một xilanh. Dấu của $A$ và $Q$ trong biểu thức của nguyên lí I Nhiệt động lực học là
	\choice
	{\True $A > 0; Q > 0$}
	{$A < 0; Q > 0$}
	{$A > 0; Q < 0$}
	{$A < 0; Q < 0$}
	\loigiai{
		
	}
\end{ex}

\begin{ex}%Câu 10
	Nhiệt dung riêng của đồng là $380\ J/kg.K$, điều này cho biết
	\choice
	{nhiệt lượng cần thiết để làm cho $1\ g$ đồng nóng lên thêm $1\ ^\circ C$ là $380\ J$}
	{nhiệt lượng cần thiết để làm cho $2\ g$ đồng nóng lên thêm $1\ ^\circ C$ là $380\ J$}
	{\True nhiệt lượng cần thiết để làm cho $1\ kg$ đồng nóng lên thêm $1\ ^\circ C$ là $380\ J$}
	{nhiệt lượng cần thiết để làm cho $1\ kg$ đồng nóng lên thêm $2\ ^\circ C$ là $380\ J$}
	\loigiai{
		
	}
\end{ex}

\begin{ex}%Câu 11
	Nhiệt độ nóng chảy của vật rắn kết tinh phụ thuộc vào
	\choice
	{\True bản chất của vật rắn và áp suất ngoài}
	{bản chất của vật rắn}
	{bản chất và nhiệt độ của vật rắn}
	{bản chất và nhiệt độ của vật rắn, đồng thời phụ thuộc áp suất ngoài}
	\loigiai{
		
	}
\end{ex}

\begin{ex}%Câu 12
	Cần đo đại lượng nào để xác định nhiệt hóa hơi riêng của nước?
	\choice
	{Khối lượng và thể tích của khối chất lỏng}
	{Nhiệt lượng cần cung cấp cho nước hóa hơi và khối lượng của nước}
	{\True Nhiệt lượng cung cấp cho khối chất lỏng trong khoảng thời gian}
	{Nhiệt lượng cung cấp cho khối chất lỏng và khoảng thời gian cung cấp nhiệt lượng đó}
	\loigiai{
		
	}
\end{ex}

\begin{ex}%Câu 13
	Phát biểu nào sau đây là \textbf{sai}. Để đảm bảo an toàn trong tiến hành thí nghiệm xác định nhiệt hoá hơi riêng của nước cần chú ý
	\choice
	{cẩn thận khi sử dụng nước nóng và nguồn điện}
	{cẩn thận khi bật tắt nguồn điện và dây điện trở}
	{chú ý quan sát mọi người xung quanh khi thao tác thí nghiệm}
	{\True nước không quá nóng và dòng điện nhỏ nên không cần chú ý}
	\loigiai{
		
	}
\end{ex}

\begin{ex}%Câu 14
	Trong thí nghiệm đo nhiệt hóa hơi riêng của nước, nhiệt hóa hơi riêng của nước được xác định bằng công thức: $L = \dfrac{\mathcal{P}.\tau}{\Delta m}$. Giá trị $\Delta m$ là
	\choice
	{khối lượng nước ban đầu trong ấm đun}
	{khối lượng nước trong ấm đun tại thời điểm $\tau$}
	{\True khối lượng nước bị bay hơi sau thời gian $\tau$}
	{khối lượng nước và ấm đun}
	\loigiai{
		
	}
\end{ex}

\begin{ex}%Câu 15
	Một đầu máy diesel xe lửa có công suất $3.10^6\ W$ và có hiệu suất là $25\%$. Cho biết năng suất tỏa nhiệt của nhiên liệu là $4,2.10^7\ J/kg$. Nếu đầu máy chạy hết công suất thì khối lượng nhiên liệu tiêu thụ trong mỗi giờ gần giá trị nào nhất sau đây?
	\choice
	{$2489\ kg$}
	{$1429\ kg$}
	{$1028\ kg$}
	{$1056\ kg$}
	\loigiai{
		\begin{itemize}
			\item $ Q_{ci} = Q_{tp}.\mathcal{H} \Rightarrow  \mathcal{P}.t= m.q. \mathcal{H} \Rightarrow 3.10^6.3600=m.4,2.10^7.0,25 \Rightarrow m \approx 1028,6\ kg$.
		\end{itemize}
	}
\end{ex}

\begin{ex}%Câu 16
	Một ấm nhôm có khối lượng $m_b = 600\ g$ chứa $V = 1,5$ lít nước ở $t_1 = 20^\circ C$, sau đó đun bằng bếp điện. Sau $t = 35$ phút thì đã có 20\% khối lượng nước đã hóa hơi ở nhiệt độ sôi $t_2 = 100^\circ C$. Biết rằng, $75\%$ nhiệt lượng mà bếp cung cấp được dùng vào việc đun nước. Cho biết nhiệt dung riêng của nước là $c_n = 4190\ J/kg.K$, của nhôm là $c_b = 880\ J/kg.K$, nhiệt hóa hơi riêng của nước ở $100^\circ C$ là $L = 2,26.10^6\ J/kg$, khối lượng riêng của nước là $D = 1$ kg/lít. Công suất cung cấp nhiệt của bếp điện gần giá trị nào nhất sau đây?
	\choice
	{\True $776\ W$}
	{$796\ W$}
	{$786\ W$}
	{$876\ W$}
	\loigiai{
		\begin{itemize}
			\item Khối lượng nước: $m_n = DV = 1,5.1 = 1,5\ kg$.
			\item Nhiệt lượng có ích: $Q = Q_n + Q_b + Q_{hh} = (m_n.c_n.+m_b.c_b)\Delta t + m_{hh}.L = (1,5.4190+0,6.880)80 + 0,2.1,5.2,26.10^6 = 1223040\ J$.
			\item Công suất: $\mathcal{P} = \dfrac{A}{t} = \dfrac{Q}{\mathcal{H}.t} = \dfrac{1223040}{0,75.35.60} \approx 776,5\ W$.
		\end{itemize}
	}
\end{ex}

\begin{ex}%Câu 17
	Bỏ $100\ g$ nước đá ở $t_1=0^\circ C$ vào $300\ g$ nước ở $t_2=20^\circ C$. Cho nhiệt nóng chảy riêng của nước đá là $\lambda=3,4.10^5\ J/kg$ và nhiệt dung riêng của nước là $c=4200\ J/kg.K$. Khối lượng đá còn lại	 bằng
	\choice
	{$0\ g$}
	{$15\ g$}
	{$21\ g$}
	{\True $26\ g$}
	\loigiai{
		\begin{itemize}
			\item Phương trình cân bằng nhiệt: $Q_{\text {tỏa }}=Q_{\text {thu }}$
			\item Khối lượng nước đá tan khi đạt cân bằng nhiệt là
			\item $m_{\tan }=\dfrac{m.c.\Delta t}{\lambda}=\dfrac{300.4200 .20}{3,4.10^5}=74\ g$
			\item Khối lượng nước đá còn lại là $m^{\prime}=m_{d}-m_{\tan }=100-74=26\ g$
		\end{itemize}
	}
\end{ex}
\begin{ex}%Câu 18
	Một ấm đun nước có công suất $500\ W$ chứa $300\ g$ nước ở $20^\circ C$. Biết nhiệt dung riêng của nước là $c = 4,18.10^3\ J/kg.K$, nhiệt hóa hơi riêng của nước $L = 2,26.10^6\ J/kg$. Nếu ấm chỉ đun trong $20$ phút. Coi như toàn bộ nhiệt lượng ấm đun được nước hấp thục hoàn toàn. Khối lượng nước còn lại trong ấm là
	\choice
	{$78,9\ g$}
	{$20,2\ g$}
	{$33,0\ g$}
	{$25,6\ g$}
	\loigiai{
		\begin{itemize}
			\item $\mathcal{P}.\tau = m.c.\Delta t + (0,3-m).L \Rightarrow 500.20.60 = 0,3.4180.80+ (0,3-m).2,26.10^6 \Rightarrow m \approx 0,0789\ kg = 78,9\ g$.
		\end{itemize}
		
	}
\end{ex}
\subsubsection*{PHẦN 2. Thí sinh trả lời từ câu 1 đến câu 4. Trong mỗi ý a), b), c), d) ở mỗi câu, thí sinh chọn đúng hoặc sai.}
\setcounter{ex}{0}
\begin{ex}
	\immini[thm]{
		Hình bên mô tả thí nghiệm làm tan chảy đá bằng nguồn nóng để xác định nhiệt nóng chảy riêng của nước đá (kết quả nhiệt nóng chảy riêng ghi lại tới chữ số thập phân đầu tiên). Biết công suất điện của nguồn nóng là $17\,W$.
		\begin{itemize}
			\item Bước 1: Ban đầu, chưa bật nguồn điện, sau $5$ phút, nước đá tan ra và chảy xuống cốc. Số chỉ hiện trên cân là $2\,g$.
			\item Bước 2: Sau đó, bật nguồn điện để nguồn nóng làm tan đá. Sau $5$ phút tiếp theo, nước đã tan ra và chảy xuống cốc. Số chỉ hiện trên cân là $20\,g$.
		\end{itemize}
	}
	{
		\begin{tikzpicture}[draw=Mapcolor, very thick, scale = 1, line join=round, line cap=round, >=stealth,line width=1]
			%\draw [dotted] (0,0) grid (6,4); 
			\draw[rounded corners=3] (0,0) -- (0.5,1) -- (2.5,1) -- (3,0)--cycle;
			\draw (0.7,1) rectangle (2.3,1.2);
			\node [rectangle,draw,minimum size=0.6cm] at (1.5,0.5) {$20.0\ g$};
			\draw[fill=cyan!30,cyan!30]  (0.5,2.5) rectangle (2.5,1.25);
			\draw[fill=cyan!5,cyan!5] (0.7,5.3) -- (2.3,5.3) -- (1.8,4.8) -- (1.2,4.8) -- (0.7,5.2);
			\draw (0.25,3) -- (0.5,2.75) -- (0.5,1.25) -- (2.5,1.25) -- (2.5,2.75) -- (2.75,3);
			\draw (0.4,5.6) -- (1.25,4.75) -- (1.35,3.5);
			
			\draw (2.6,5.6) -- (1.75,4.75) -- (1.65,3.5);
			\draw [fill=cyan!30,cyan!30] (1.5,3.25) .. controls (1.25,2.75) and (1.75,2.75) .. (1.5,3.25);
			\draw[fill=cyan!40]  (0.8,5.3) .. controls (0.7,5.4) and (1.1,5.4) .. (1.1,5.2) .. controls (1,5) and (0.8,5.2) .. (0.8,5.3);
			\draw[fill=cyan!40]  plot[smooth, tension=.7] coordinates {(1.3,5) (1.2,5.2) (1.4,5.2) (1.6,5.1) (1.5,4.9) (1.3,5)};
			\draw [fill=cyan!40] plot[smooth, tension=.7] coordinates {(2,5.2) (1.8,5.4) (1.7,5.3) (1.6,5) (2,5) (2,5.2)};
			\draw [fill=cyan!40] plot[smooth, tension=.7] coordinates {(1.3,5.1) (1,5) (1.2,4.9) (1.4,5) (1.3,5.1)};
			\draw [fill=cyan!40] plot[smooth, tension=.7] coordinates {(1.3,5.5) (1.1,5.3) (1.3,5.2) (1.5,5.3) (1.4,5.5) (1.3,5.5)};
			\draw [fill=cyan!40] plot[smooth, tension=.7] coordinates {(1.7,5.6) (1.5,5.5) (1.6,5.3) (2,5.4) (1.8,5.7) (1.7,5.6)};
			\draw [fill=cyan!40] plot[smooth, tension=.7] coordinates {(2.4,5.4) (2.3,5.6) (2,5.6) (2.1,5.4) (2.2,5.3) (2.4,5.4)};
			\draw[line width=1.5,gray] (0,6.5) .. controls (1.5,6.5) and (1.5,6) .. (1.5,5);
			\draw (1.44,5.8) -- (3.8,6.4);
			\draw (1.8,5.2) -- (3.8,4.6);
			\draw (2.5,2.2) -- (4.2,2.2);
			\draw (2.7,0.6) -- (4.4,0.4);
			\node at (5,6.4) {Nguồn nhiệt};
			\node at (4.8,4.6) {Nước đá};
			\node at (5.2,2.2) {Bình chứa};
			\node at (5.4,0.4) {Cân điện tử};
		\end{tikzpicture}    	
	}
	\choiceTF[t]
	{Trong $5$ phút đầu tiên (ở bước $1$), do chưa bật nguồn điện nên nội năng của nước đá không đổi}
	{Sau khi bật nguồn, khối lượng của nước đá đã tan là $20\,g$}
	{Nhiệt nóng chảy riêng của nước đá tính ra trong thí nghiệm trên là $318,75\,kJ/(kg.K)$}
	{\True Biết giá trị nhiệt nóng chảy riêng của nước đá được ghi nhận là $333\,kJ/kg$. Sai lệch của kết quả thí nghiệm so với giá trị đã được ghi nhận là $4,28\%$}
	\loigiai{
		\begin{itemchoice}
			\itemch Trong $5$ phút đầu tiên (ở bước $1$), do chưa bật nguồn điện nhưng nước đá hấp thụ nhiệt từ môi trường, nội năng của nước đá vẫn tăng.
			\itemch Sau khi bật nguồn, số chỉ hiện trên cân là $20\,g$ là tổng khối lượng nước đá tan do nhận nhiệt từ nguồn điện và nhiệt từ môi trường cộng với lượng nước đá tan trước khi bật nguồn điện.
			\itemch Lượng nước đã thực tế đã tan do nhận nhiệt từ nguồn nóng sau $5$ phút bật nguồn điện hay $10$ phút tính từ thời điểm ban đầu là $20 - 2 - 2 = 16\,g.$ Nhiệt nóng chảy riêng của nước đá tính ra trong thí nghiệm trên là $\lambda = \dfrac{Q}{m} = \dfrac{17 . 5 . 60}{16 . 10^{-3}} = 318,75\,kJ/kg \Rightarrow$ sai đơn vị. 
			\itemch Sai lệch của kết quả thí nghiệm so với giá trị đã được ghi nhận là $\dfrac{318,75}{333} . 100\% = 4,28\%$.
		\end{itemchoice}
	}
	
\end{ex}

\begin{ex}
	Người ta nung nóng miếng đồng có khối lượng $100\ g$ đến nhiệt độ $650^\circ C$ rồi thả vào cốc nước có thể tích 1 lít đang có nhiệt độ $30^\circ C$. Giả sử cốc nước được làm từ vật liệu cách nhiệt và bỏ qua sự trao đổi nhiệt giữa khối nước và môi trường bên ngoài. Biết khối lượng riêng của nước là $1000\ \text{kg/m}^3$ và nhiệt dung riêng của đồng và nước lần lượt là $c_1 = 380\ J/kg.K$ và $c_2 = 4180\ J/kg.K$.
	\choiceTF[t]
	{\True Đã có quá trình truyền nhiệt từ miếng đồng sang nước}
	{\True Khi hệ đã cân bằng nhiệt, so với ban đầu thì nội năng của miếng đồng đã giảm xuống, còn của nước tăng lên}
	{\True Khi hệ đã cân bằng nhiệt, nhiệt độ của nước trong cốc là $35,59^\circ C$}
	{Độ biến thiên nội năng của miếng đồng là $24700\ J$}
	\loigiai{
		\begin{itemchoice}
			\itemch Đúng. Miếng đồng có nhiệt độ cao hơn nên truyền nhiệt cho nước.
			\itemch Đúng. Đồng mất nhiệt (nội năng giảm), nước nhận nhiệt (nội năng tăng).
			\itemch Đúng. Giải phương trình cân bằng nhiệt ta được $t_{cb} \approx 35,59^\circ C$.
			\itemch Sai. $\Delta U = m.c.\Delta t = 0,1.380.(650 - 35,59) \approx 23348\ J$
		\end{itemchoice}
	}
\end{ex}

\begin{ex}
	Một nhóm học sinh thực hiện thí nghiệm do nhiệt dung riêng nóng chảy của nước đá và lập được bảng số liệu như sau:
	
	\begin{center}
		\begin{tabular}{|l|c|c|c|c|c|c|c|c|c|}
			\hline
			Thời gian $t(s)$ & 0 & 120 & 240 & 360 & 480 & 600 & 720 & 840 & 960 \\
			\hline
			Nhiệt độ $t \left({ }^\circ C\right)$ & 0 & 0 & 0 & 0 & 0 & 0 & 0,3 & 0,6 & 0,9 \\
			\hline
			Công suất $P(W)$ & 14,25 & 14,25 & 14,25 & 14,25 & 14,25 & 14,25 & 14,25 & 14,25 & 14,25 \\
			\hline
		\end{tabular}
	\end{center}
	Biết khối lượng nước đá $m=0,025\ kg$ và bỏ qua sự truyền nhiệt giữa nhiệt lượng kế với môi trường.
	\choiceTF[t]
	{\True Nhiệt lượng đā cung cấp cho nước đá trong 120 s đầu tiên là $1710\ J$}
	{\True Thời gian để nước đá nóng chảy hoàn toàn là 600 s}
	{Dựa vào bảng số liệu, nhiệt nóng chảy riêng của nước đá là $337000\ J/kg$}
	{\True Dựa vào nhiệt nóng chảy của nước đá tính được trong thí nghiệm trên. Nếu người ta cung cấp nhiệt lượng $10,02.10^5 J$ và hiệu suất quá trình $80 \%$ thì khối lượng nước đá có thể làm nóng chảy hoàn toàn là $2,34\ kg$}
	\loigiai{
		\begin{itemchoice}
			\itemch Nhiệt lượng đā cung cấp cho nước đá trong 120 s đầu tiên là $Q=P . t=14,25.120=1710(J)$
			\itemch Thời gian để nước đá nóng chảy hoàn toàn là 600 s.
			\itemch Dựa vào bảng số liệu, nhiệt nóng chảy riêng của nước đá thỏa mān $Q=P. t=m . \lambda \Rightarrow \lambda=\dfrac{P. t}{m}=\dfrac{600.14,25}{0,025}=342000\ J/kg$
			\itemch Dựa vào nhiệt nóng chảy của nước đá tính được trong thí nghiệm trên. Nếu người ta cung cấp nhiệt lượng $10,02.10^5\ J$ và hiệu suất quá trình $80 \%$ thì khối lượng nước đá có thể làm nóng chảy hoàn toàn là $Q=m . \lambda \Rightarrow m=\dfrac{Q}{\lambda}=\dfrac{10,02.10^5 .80 \%}{342000}=2,34\ kg$
		\end{itemchoice}
	}
\end{ex}

\begin{ex}
	\immini[thm]{Một nhóm học sinh làm thí nghiệm để xác định nhiệt dung riêng của một mẫu kim loại. Họ có một bình xốp hình trụ có vỏ và nắp cách nhiệt, một que khuấy, một nhiệt kế, mẫu kim loại, một chiếc cân và một bình đun nước. Ban đầu, mẫu kim loại được để ở nhiệt độ phòng ($27,0^\circ C$).}{\begin{tikzpicture}[ thick]
			%\draw [dotted] (-2,0) grid (2,5);
			\foreach \x in {0.01,0.02,...,0.1}{
				\fill [opacity=0.4] (0,1.5+\x) ellipse (1.15cm and 0.4cm);}
			\foreach \x in {0.1,0.11,...,0.2}{
				\fill [ball color = gray!50, scale = 0.5] (0,\x) ellipse (1.25cm and 0.5cm);}
			\fill [rounded corners=2pt, opacity=0.7] (-1,1.5) rectangle (-0.8,0.5);
			\fill [ball color = white,rounded corners = 2pt] (-0.1,0.5) rectangle (0.1,3.3);
			\fill [ball color = red] (0,0.4) circle (0.2cm);
			\fill [ball color = red] (-0.05,0.6) rectangle (0.05,2.3);
			\draw
			(0,3) ellipse (1.25cm and 0.5cm)
			(0,3) ellipse (1.15cm and 0.4cm)
			(-1.25,3) -- (-1.25,0) arc [start angle = -180, end angle = 0, x radius = 1.25, y radius = 0.5] (1.25,0) -- (1.25,3);
			\draw [dashed]
			(-1.25,0) arc [start angle = 180, end angle = 0, x radius = 1.25, y radius = .5]
			(-1.25,1.5) arc [start angle = 180, end angle = 0, x radius = 1.25, y radius = .5];
			\fill [blue!50, opacity=0.5] (-1.15,1.5) -- (-1.15,0) arc [start angle = -180, end angle = 0, x radius = 1.15, y radius = 0.4] (1.15,0) -- (1.15,1.5) arc [start angle = 0, end angle = -180, x radius = 1.15, y radius = 0.4];
			\fill [blue!50, opacity=0.5] (0,1.5) ellipse (1.15cm and 0.4cm);
			\fill [ opacity=0.8] (-1,3.5) rectangle (-0.8,1.5);
			\draw (-1.25,1.5) arc[start angle = -180, end angle = 0, x radius = 1.25, y radius = 0.5];
			\draw [dashed] (-1.15,3) -- (-1.15,0) arc [start angle = -180, end angle = 0, x radius = 1.15, y radius = 0.4] (1.15,0) -- (1.15,3);
			\draw
			(-0.9,3.5) -- (-2,4) node[above] {que khuấy}
			node[above] at (0,3.5) {nhiệt kế}
			(1.1,1.5) -- (2,1.5) node[right, text width = 2.5 cm, align = center] {vỏ và nắp cách nhiệt}
			(0,0) -- (2,0) node[right] {mẫu kim loại};
			
	\end{tikzpicture}}
	\choiceTFt
	{Nhóm học sinh sử dụng cân và xác định được khối lượng nước đổ vào bình xốp là $0,225 \ kg$, khối lượng của mẫu kim loại là $0,409 \ kg$. Số chỉ của nhiệt kế nhúng trong nước nóng ngay trước khi thả mẫu kim loại là $67,5^\circ C$ và số chỉ của nhiệt kế khi mẩu kim loại và nước đạt trạng thái cân bằng nhiệt là $56, 0^\circ C$. Biết nhiệt dung riêng của nước là $4180 \ J/kg.K$. Từ các số liệu trên, nhóm học sinh xác định được nhiệt dung riêng của mẫu kim loại là $889 \ J/kg.K$}
	{\True Nhóm học sinh cho rằng, nếu đun nóng nước tới khoảng $70, 0^\circ C$, đổ vào bình xốp đã cắm sẵn nhiệt kế, nhẹ nhàng nhúng chìm mẫu kim loại trong nước, đóng kín nắp lại và khuấy nhẹ tay thì số chỉ trên nhiệt kế sau đó sẽ thay đổi liên tục và chỉ dừng lại khi bình xốp chứa nước cùng mẫu kim loại đạt trạng thái cân bằng nhiệt}
	{Nhóm học sinh cho rằng, kết quả tính được ở câu a) nhỏ hơn giá trị nhiệt dung riêng chính xác của mẫu kim loại do trong phép tính đã bỏ qua nhiệt lượng trao đổi với môi trường}
	{Một học sinh trong nhóm cho rằng, nếu bỏ qua thất thoát nhiệt với môi trường thì nhiệt lượng nước thu vào bằng với nhiệt lượng mẫu kim loại tỏa ra}
	\loigiai{
		\begin{itemchoice}
			\itemch $m_{kl}.c_{kl}.\Delta t_{kl} = m_n.c_n.\Delta t_n \Rightarrow 0,409.c_{kl}.29 = 0,225.4180.11,5 \Rightarrow c_{kl} \approx 912\ J/kg.K$.
			\itemch $Q = m.q = 1,5.4,6.10^7 = 69.10^6\ J = 6,9.10^7\ J$.
			\itemch $0,409.c_{kl}.29 + Q_{mt} = 0,225.4180.11,5 \Rightarrow c_{kl} = \dfrac{0,225.4180.11,5 - Q_{mt}}{0,409.29} \Rightarrow$ kết quả tính được ở câu a) lớn hơn giá trị nhiệt dung riêng chính xác mẫu kim loại.
			\itemch Nếu bỏ qua thất thoát nhiệt ra môi trường thì nhiệt lượng nước tỏa ra bằng nhiệt lượng mẫu kim loại thu vào.
		\end{itemchoice}
	}
\end{ex}
\subsubsection*{PHẦN 3. Thí sinh trả lời từ câu 1 đến câu 6.}
\setcounter{ex}{0}
\begin{ex}% Câu 1
	Một khối khí được truyền một nhiệt lượng $2000\ J$ thì khối khí dãn nở và thực hiện được một công $1500\ J$. Độ biến thiên nội năng của khối khí là bao nhiêu $J$ (làm tròn kết quả đến chữ số hàng đơn vị)?
	\shortans[oly]{$500$}
	\loigiai{
		\begin{itemize}
			\item Khối khí được truyền nhiệt lượng nên $Q = 2000\ J$.
			\item Khối khí dãn nở và thực hiện công nên $A = -1500\ J$.
			\item Độ biến thiên nội năng của khối khí là $\Delta U = Q + A = 2000 + (-1500) = 500\ (J)$.
		\end{itemize}
	}
\end{ex}

\begin{ex}% Câu 2
	Có hai bình cách nhiệt. Bình I chứa $5\ l$ nước ở $60^\circ C$, bình II chứa $1\ l$ nước ở $20^\circ C$. Đầu tiên, rót một phần nước ở bình I sang bình II. Sau khi bình II cân bằng nhiệt, người ta lại rót từ bình II sang bình I một lượng nước bằng với lần rót trước. Nhiệt độ sau cùng của nước trong bình I là $59^\circ C$. Thể tích nước đã rót từ bình này sang bình kia bằng bao nhiêu lít (làm tròn kết quả đến chữ số hàng phần trăm)?
	\SA[oly]{$0,14$}
	\loigiai{
		Gọi $x$ (lít) là lượng nước đã rót từ bình này sang bình kia.
		\begin{itemize}
			\item Cân bằng nhiệt ở bình II: $x.c(60-t) = 1.c(t-20) \Rightarrow t = \dfrac{60x+20}{x=1}\ (1)$.
			\item Cân bằng nhiệt ở bình I: $(5-x).c.1 = x.c(59-t) \Rightarrow t= \dfrac{60x-5}{x}\ (2)$.
			\item Từ (1) và (2) $x = \dfrac{1}{7} \approx 0,14$ lít.
		\end{itemize}
	}
\end{ex}

\begin{ex}% Câu 3
	\immini[thm]{Hai chất có cùng nhiệt dung riêng $c$ với khối lượng tương ứng là $m_1$ và $m_2$ được cung cấp nhiệt với công suất không đổi $\mathcal{P}$. Đồ thị bên mô tả sự thay đổi nhiệt độ của chúng theo thời gian $t$ cung cấp nhiệt. Tỷ lệ $\dfrac{m_1}{m_2}$ có giá trị (làm tròn kết quả đến chữ số hàng phần mười)?}{\begin{tikzpicture}[draw=Mapcolor, very thick, >=latex']
			\draw [dotted] (0,0) grid (5,5);
			\draw [->] (0,0) -- (0,5.5) node [right] {$\theta$};
			\draw [->] (0,0) -- (5.5,0) node [right] {$\tau$};
			\coordinate (O) at (0,0);
			\coordinate (A) at (3,4);
			\coordinate (B) at (5,3);
			\coordinate (M) at (5,0);
			\draw [line width=1.5pt] (O) -- (A) node [above] {$m_2$}
			(O) -- (B) node [above] {$m_1$};
			\draw pic[ draw,-, angle radius=1cm, angle eccentricity=1.2] {angle = M--O--A};
			\draw pic[ draw,-, angle radius=1.1cm, angle eccentricity=1.2] {angle = M--O--A};
			\draw pic[ draw,-, angle radius=0.5cm, angle eccentricity=1.2] {angle = M--O--B};
			\draw 
			node [right] at (15:0.5) {$\alpha _1$}
			node [right] at (26:1.1) {$\alpha _2$};
	\end{tikzpicture}}
	\SA[oly]{$2,2$}
	\loigiai{
		\begin{itemize}
			\item $\mathcal{P}.t = m.c.\theta \Rightarrow \dfrac{t_1}{t_2}=\dfrac{m_1}{m_2}.\dfrac{\theta _1}{\theta _2} \Rightarrow \dfrac{5}{3} = \dfrac{m_1}{m_2}.\dfrac{3}{4} \Rightarrow \dfrac{m_1}{m_2} = \dfrac{20}{9} \approx 2,2$.
		\end{itemize}
	}
\end{ex}

\begin{ex}% Câu 4
	Người ta đổ $m_1 = 200\ g$ nước sôi có nhiệt độ $100^\circ C$ vào một chiếc cốc thủy tinh có khối lượng $m_2 = 120\ g$ đang ở nhiệt độ $t_2 = 20^\circ C$. Sau khoảng thời gian $t = 5$ phút, nhiệt độ của cốc nước bằng $40^\circ C$. Nhiệt dung riêng của nước là $c_1 = 4200\ J/kg.K$ và của thuỷ tinh là $c_2 = 840\ J/kg.K$. Xem rằng sự mất mát nhiệt xảy ra một cách đều đặn, nhiệt lượng toả ra môi trường xung quanh trong mỗi giây bằng bao nhiêu $J$ (làm tròn kết quả đến chữ số hàng đơn vị)?
	\SA[oly]{$161$}
	\loigiai{
		\begin{itemize}
			\item $Q_n = Q_{tt} + Q_{hp} \Rightarrow m_n.c_n.\Delta t_n = m_{tt}.c_{tt}.\Delta t_{tt} + \mathcal{P}_{mt}.\tau \Rightarrow 0,2.4200.60 = 0,12.840.20 +\mathcal{P}.5.60 \Rightarrow \mathcal{P} = 161,26\ W$.
		\end{itemize}
	}
\end{ex}

\begin{ex}% Câu 5
	Vào mùa hè, người Hà Nội thường có thói quen thưởng thức trà đá trong các quán vỉa hè. Để có một cốc trà đá chất lượng, người chủ quán rót khoảng $0,250\ kg$ trà nóng ở $80,0\ ^{\circ}C$ vào cốc, sau đó cho tiếp $m\ kg$ nước đá $0\ ^{\circ}C$. Cuối cùng được cốc trà đá ở nhiệt độ phù hợp nhất là $10,0\ ^{\circ}C$ (hệ vừa đạt đến trạng thái cân bằng nhiệt). Biết phần nhiệt lượng mà hệ (nước và nước đá) nhận thêm của môi trường xung quanh bằng $10\%$ nhiệt lượng mà các cục nước đá nhận để làm tăng nội năng của chúng. Nhiệt dung riêng của nước là $4,20\ kJ/kg.^{\circ}C$; nhiệt nóng chảy của nước đá là $3,33.10^5\ J/kg$. Giá trị $m$ bằng bao nhiêu $kg$ (làm tròn kết quả đến chữ số hàng phần trăm)?
	\SA[oly]{$0,22$}
	\loigiai{
		
		\begin{itemize}
			\item Nhiệt lượng nước tỏa ra + Nhiệt lượng môi trường cung cấp = Nhiệt lượng nước đá nhận vào để nóng chảy và tăng nhiệt độ:
			\begin{align*}
				&Q_n +0,1\left( Q_{nc} + Q_\text{đ}\right) =  \left( Q_{nc} + Q_\text{đ}\right) \Rightarrow Q_n = 0,9 \left( Q_{nc} + Q_\text{đ}\right)\\
				\Rightarrow &m_n.c_n.\Delta t_n = 0,9m\left( \lambda + c_n.t \right)\\
				\Rightarrow &0,25.4,2.10^3.70 = 0,9.m(3,33.10^5+4,2.10^3.10) \Rightarrow m \approx 0,22\ kg
			\end{align*}
			
		\end{itemize}
	}
\end{ex}

\begin{ex}% Câu 6
	Một bình cổ cong đựng nước ở $0\ ^\circ\mathrm{C}$. Người ta làm nước trong bình đông đặc lại bằng cách hút không khí và hơi nước trong bình ra ngoài. Lấy nhiệt nóng chảy riêng của nước là $3,3. 10^5\ J/kg$ và nhiệt hóa hơi riêng ở nước là $2, 48.10^6\ J/kg$. Phần nhiệt lượng trao đổi nhiệt với môi trường bên ngoài (nhiệt độ môi trường bên ngoài cao hơn $0\ ^\circ\mathrm{C}$) thông qua thành bình chiếm $20\%$ năng lượng mà lượng nước đông đặc tỏa ra. Tỉ số giữa khối lượng nước bị hóa hơi tối đa và khối lượng nước ở trong bình lúc đầu là bao nhiêu (làm tròn kết quả đến chữ số hàng phần trăm)?
	\SA[oly]{$0,14$}
	\loigiai{
		\begin{itemize}
			\item Nhiệt lượng nước thu vào để hóa hơi: $Q_h = m_h.L$.
			\item Nhiệt lượng nước tỏa ra để đông đặc: $Q_\text{đđ} = \lambda (m-m_h)$.
			\item Phương trình cân bằng nhiệt: $Q_h = Q_\text{đđ}+0,2Q_\text{đđ} \Rightarrow m_h.L = 1,2\lambda (m-m_h) \Rightarrow 2,48.10^6.m_h = 1,2.3,3.10^5 (h-m_h) \Rightarrow \dfrac{m_h}{m} \approx 0,14$.
		\end{itemize}
	}
\end{ex}
